\documentclass[11pt]{article}
\usepackage[margin=0.85in]{geometry}
\usepackage[T1]{fontenc}
\usepackage{array}
\usepackage{amsmath}
\usepackage{amssymb}
\setlength{\parskip}{5pt}
\setlength{\parindent}{0pt}
\begin{document}
\begin{center}
{\Large Explicit infinity point: checked proof boundary}\\[4pt]
{\small aburan28/crypto, genus-2 projective-law follow-up, 25 October 2026}
\end{center}

\textbf{Object.} For a field $K$ and arbitrary $a_3,a_2,a_1,a_0\in K$, the
explicit infinity chart has ring
\[
  R=K[z,w]/\bigl(w^2-z(1+a_3z^2+a_2z^3+a_1z^4+a_0z^5)\bigr).
\]
The checked overlap with the affine chart is the principal open $D(z)$.

\textbf{Checked result.} The complement of $D(z)$ in $\operatorname{Spec}R$
contains exactly one point, the origin ideal $(z,w)$. For any prime $\mathfrak p$
containing $z$, the equation gives $w^2\in\mathfrak p$, and primality gives
$w\in\mathfrak p$. The pre-existing Lean theorem proves $(z,w)$ is maximal,
so $\mathfrak p=(z,w)$. Evaluation at $z=w=0$ proves the origin is outside
$D(z)$. This statement is for the explicit affine chart; it is not a claim
of projectivity or smoothness of the glued scheme.

\textbf{Compiler evidence.} Pinned Lean v4.34.1 and Mathlib revision
\begin{center}\small\texttt{d13f23b723b8a846827a245b89c10fc7d3f11612}\end{center}
The command
\texttt{timeout 180s lake env lean Common22CompletionBridge.lean} exited 0.
The four \texttt{\#print axioms} lines list exactly
\texttt{propext}, \texttt{Classical.choice}, and \texttt{Quot.sound}.
The adjacent CSV records each theorem, its hypotheses, and status; the SVG
shows the dependency and open-obligation graph.

\textbf{Remaining bridge.} A chosen smooth projective completion must be
constructed and its two affine charts identified with the present quotient
rings, with the same overlap isomorphism. Alternatively, construct a finite
map from the glued scheme to $\mathbf P^1_K$, then prove the properness and
chartwise smoothness hypotheses under an explicit squarefree condition.
The pinned Mathlib supplies generic gluing and $\operatorname{Proj}$ APIs,
but no genus-2 hyperelliptic completion theorem was found. The principal
divisor of the interpolating function, including its order at this unique
point, and the resulting $\operatorname{Pic}^0$ equation remain separate
proof obligations.

\textbf{Interpretation.} This packet advances the local geometry of the
constructed chart. It does not revise measured cost ratios or assert
correctness of the full published law.
\end{document}
